\scanpage{22}\mspage{18 bis}
% @ top, opening | text | Soit r un entier
Soit $r$ un entier,

% @ top, a | text | Q un ens.
a) $Q$ un ens. de cardinal $r$

% @ top, b | text | Q tilde, π, Q
b) $\widetilde Q\xrightarrow{\pi} Q$ un rev. d'ordre 2 de $Q$, $\Longleftrightarrow$ un torseur \deleted{$Q$} \inserted{$\Pi$} sous $\mathbb F_2^Q$

% @ upper third, crossed c and its continuation | text | restriction du groupe structural
\deleted{c) restriction du groupe structural à $\Psi_Q=\operatorname{Ker}\bigl(\mathbb F_2^Q\longrightarrow\mathbb F_2\bigr)$, $(x_i)\longmapsto\sum x_i$,

La donnée de b), c) moyennant a), revient à la donnée d'un torseur $P$ sous $\Psi_Q$.

On va associer à ceci un ens de racines de type $D_r$.}

% @ upper middle, new a | text | produit scalaire
a) Soit $(x,y)$ le produit scalaire sur $\mathbb Z^Q$ défini pour $x,y\in Q$ par
% @ upper middle, cases | display | x.y = 0, 1, −1
\[
x.y=\begin{cases}0&\text{si }\deleted{\unsure{y}}\quad y\ne x,\sigma x\ \text{i.e. }\pi(x)\ne\pi(y)\\
1&\text{si }\quad y=\sigma x\\
-1&\text{si }\quad y=x\end{cases}
\]
% @ centre, quotient | text | Soit E(Π)
Soit $E(\unsure{\Pi|P|\deleted{\Pi}})=\mathbb Z^Q/\{x\in\mathbb Z^Q\mid x.y=0\quad\forall y\in\mathbb Z^Q\}$

% @ centre, cancelled choices | text | Choisissons ; Ordonnons Q
\deleted{Choisissons un élément} \deleted{Ordonnons $Q\simeq[1,r]$ et choisissons un élément $\tilde q_{\unsure{i}}\in\widetilde Q$ sur tt $q\in Q$, \deleted{alors on aura} soit $\sigma\tilde q$ lié à $\tilde q$, $E$ $\tilde q_i^2=-1\quad\tilde q_i\tilde q_j=0$} je dis que dans $E$, $\tilde q+\sigma\tilde q=0$ \deleted{indépendant de $q$, car} car le produit scalaire \deleted{$(\sigma\tilde q+\tilde q,\tilde q')=\bigl\{$} avec $\tilde q$ est $\deleted{\unsure{+}}\unsure{-|{}}1+1$, avec $\sigma\tilde q$ est $+1-1$, et avec \deleted{\unsure{a}} $\tilde q'\ne\tilde q,\sigma\tilde q$ est $0+0$. Donc $E$ est le $\mathbb Z$-module libre de rang $r$ engendré par les $\tilde q\in\widetilde Q$, avec $\sigma\tilde q=-\tilde q$, $(\tilde q,\tilde q)=-1$, $(\tilde q,\tilde q')=0$ si $\tilde q'\ne\tilde q,\sigma\tilde q$ \ldots.

% @ lower third, roots | text | Les racines sont les différences
Les racines sont les différences mutuelles des $\tilde q-\tilde q'$, avec $\tilde q'\ne\tilde q,\sigma\tilde q=-\tilde q$

% @ foot, cubic complex | text | Supposons que Q
Supposons que $Q=\unsure{F(u)|E(u)}$, $u$ sommet d'un complexe cubique \inserted{$C$}, OPS $C$ muni d'une base $(\unsure{\zeta|\xi}_i)_{1\le i\le6}$ et \deleted{\unsure{un}} $u=\zeta'_6$, donc
% @ foot, array E(u) | display | ζ₁ through ζ₅₆
\[
E(u)=\begin{pmatrix}\zeta_1&\zeta_2&\zeta_3&\zeta_4&\zeta_5\\\zeta_{16}&\zeta_{26}&\zeta_{36}&\zeta_{46}&\zeta_{56}\end{pmatrix},
\]
% @ very bottom, generators | text | Soit E
\deleted{\unsure{c'est}} Soit $\unsure{E'|E_f}\subset\frac12 E(C)$ engendré par les \deleted{$\varepsilon_i=\frac12(\eta-2\zeta_i-\zeta_6)=\frac12(\eta-$} $\varepsilon_i=\underbrace{\frac12(\zeta_6-\eta)}_{\theta_0}+\zeta_i$, alors les diff. mutuelles dans $\widehat E(C)$ des couples d'élts non liés de $E(u)$ sont les
% archive sheet number: 21
